%═══════════════════════════════════════════════════════════════════════════════
\chapter{Geometría del espacio de parámetros y paisajes de pérdida}
\label{ch:cap3}
%═══════════════════════════════════════════════════════════════════════════════

El Capítulo~\ref{ch:cap2} cerró con una advertencia: cuando la familia paramétrica
es no lineal en $\bm{\theta}$, la pérdida empírica deja de ser una forma cuadrática
convexa. Esto significa que el paisaje de $\mathcal{L}:\mathbb{R}^p \to \mathbb{R}$
puede tener valles, crestas, mesetas y múltiples hoyos. Antes de diseñar un
algoritmo para descender por ese paisaje, conviene estudiarlo con cuidado.

Este capítulo desarrolla las herramientas del cálculo diferencial en $\mathbb{R}^p$
que describen la geometría local y global de $\mathcal{L}$. El gradiente dice en
qué dirección la función crece más rápido. El hessiano describe la curvatura local.
La convexidad es la propiedad que garantiza un paisaje sin trampas. Y la no
convexidad, que es la situación real de las redes neuronales, tiene una estructura
en alta dimensión que resulta sorprendentemente manejable por razones que solo se
ven cuando $p$ es grande.

Para el físico, este capítulo es cálculo diferencial en alta dimensión: las mismas
herramientas que se usan para analizar potenciales de energía en mecánica, campos
escalares en electrodinámica o funcionales en mecánica cuántica, aplicadas ahora
a la función de pérdida como objeto de estudio en sí mismo.

%───────────────────────────────────────────────────────────────────────────────
\section{La pérdida empírica como función escalar en \texorpdfstring{$\mathbb{R}^p$}{Rp}}
\label{sec:loss_funcion}
%───────────────────────────────────────────────────────────────────────────────

\dificultad{1}{1}

\begin{definicion}{Pérdida empírica como función de los parámetros}
\label{def:loss_funcion}
Sea $f:\mathbb{R}^n \times \mathbb{R}^p \to \mathbb{R}$ una familia paramétrica
y $\mathcal{D} = \{(\bm{x}_i, y_i)\}_{i=1}^N$ un conjunto de datos. La
\textbf{pérdida empírica} se considera como función exclusivamente de $\bm{\theta}$:
\[
  \mathcal{L}: \mathbb{R}^p \longrightarrow \mathbb{R},
  \qquad
  \mathcal{L}(\bm{\theta}) = \frac{1}{N}\sum_{i=1}^N \ell\bigl(f(\bm{x}_i;\bm{\theta}),\,y_i\bigr).
\]
Los datos $\mathcal{D}$ son fijos; solo $\bm{\theta}$ varía. El dominio es
$\mathbb{R}^p$, donde $p$ es el número de parámetros del modelo.
\end{definicion}

La diferenciabilidad de $\mathcal{L}$ se hereda de la de sus componentes por la
regla de la cadena.

\begin{proposicion}{Diferenciabilidad de la pérdida empírica}
\label{prop:diff_loss}
Si $f(\bm{x}_i;\cdot):\mathbb{R}^p \to \mathbb{R}$ es de clase $C^k$ en $\bm{\theta}$
para todo $i$, y la función de pérdida $\ell:\mathbb{R}\times\mathbb{R}\to\mathbb{R}$
es de clase $C^k$, entonces $\mathcal{L} \in C^k(\mathbb{R}^p)$.
\end{proposicion}

\begin{proof}
$\mathcal{L}$ es suma finita de composiciones $\bm{\theta} \mapsto
\ell(f(\bm{x}_i;\bm{\theta}), y_i)$. Cada composición es $C^k$ por la regla de
la cadena para funciones $C^k$. La suma finita preserva la clase de diferenciabilidad.
\end{proof}

En casi todos los modelos de interés práctico se tiene $\mathcal{L} \in C^2(\mathbb{R}^p)$:
la pérdida MSE con familia polinomial es $C^\infty$; con activaciones suaves como
tanh o sigmoide la red neuronal también lo es; solo las activaciones no diferenciables
como ReLU rompen esta regularidad en un conjunto de medida cero.

\subsection{Superficies de nivel y conjuntos de sublevel}
\label{subsec:nivel}

\dificultad{1}{2}

La geometría de $\mathcal{L}$ se visualiza a través de sus superficies de nivel.

\begin{definicion}{Superficie de nivel y conjunto de sublevel}
\label{def:nivel}
Sea $c \in \mathbb{R}$. La \textbf{superficie de nivel} de $\mathcal{L}$ para el
valor $c$ es:
\[
  \mathcal{S}_c = \bigl\{\bm{\theta} \in \mathbb{R}^p : \mathcal{L}(\bm{\theta}) = c\bigr\}.
\]
El \textbf{conjunto de sublevel} es:
\[
  \mathcal{L}_c = \bigl\{\bm{\theta} \in \mathbb{R}^p : \mathcal{L}(\bm{\theta}) \leq c\bigr\}.
\]
\end{definicion}

Para la familia lineal con pérdida MSE, $\mathcal{L}(\bm{\theta}) =
\frac{1}{N}\|\bm{V}\bm{\theta}-\bm{y}\|_2^2$ es una forma cuadrática: sus
superficies de nivel son \textbf{elipsoides} en $\mathbb{R}^p$ centrados en
$\bm{\theta}^*$, y sus ejes están alineados con los vectores singulares derechos
$\bm{w}_i$ de $\bm{V}$ con longitudes proporcionales a $1/\sigma_i$. El mínimo
se encuentra en el centro del elipsoide. Para una red neuronal, la misma geometría
local existe en torno a cada punto crítico, pero globalmente las superficies de nivel
pueden ser no conexas y tener topología complicada.

\begin{observacion}{Alta dimensión e intuición geométrica}
En $\mathbb{R}^2$ o $\mathbb{R}^3$ se puede visualizar directamente la forma de
$\mathcal{L}$. Para $p = 10^6$ la visualización directa es imposible, y la
intuición geométrica formada en baja dimensión puede ser engañosa. Un resultado
famoso de geometría en alta dimensión ilustra esta ruptura de intuición: casi toda
la volumen de una bola de radio $r$ en $\mathbb{R}^p$ está concentrada en una
cáscara delgada de grosor $O(r/p)$ cerca de la superficie. Los métodos de análisis
de alta dimensión (concentración de medida, geometría estocástica) revelan que las
redes neuronales operan en un régimen cuyas propiedades son cualitativamente
distintas de las que la intuición de baja dimensión sugiere. A lo largo del capítulo
se señalarán los puntos donde esa distinción es relevante.
\end{observacion}

%───────────────────────────────────────────────────────────────────────────────
\section{El gradiente: dirección y magnitud del cambio}
\label{sec:gradiente}
%───────────────────────────────────────────────────────────────────────────────

\dificultad{2}{2}

\begin{definicion}{Gradiente de la pérdida}
\label{def:gradiente}
Sea $\mathcal{L} \in C^1(\mathbb{R}^p)$. El \textbf{gradiente} de $\mathcal{L}$
en $\bm{\theta}$ es el vector:
\[
  \nabla_{\bm{\theta}}\mathcal{L}(\bm{\theta})
  = \left(\frac{\partial\mathcal{L}}{\partial\theta_1},\;
           \frac{\partial\mathcal{L}}{\partial\theta_2},\;
           \ldots,\;
           \frac{\partial\mathcal{L}}{\partial\theta_p}\right)^\top
  \in \mathbb{R}^p.
\]
La \textbf{derivada direccional} de $\mathcal{L}$ en $\bm{\theta}$ a lo largo
de un vector unitario $\bm{v} \in \mathbb{R}^p$ ($\|\bm{v}\|_2 = 1$) es:
\[
  D_{\bm{v}}\mathcal{L}(\bm{\theta})
  = \lim_{h \to 0} \frac{\mathcal{L}(\bm{\theta}+h\bm{v}) - \mathcal{L}(\bm{\theta})}{h}
  = \bigl\langle \nabla_{\bm{\theta}}\mathcal{L}(\bm{\theta}),\, \bm{v} \bigr\rangle.
\]
\end{definicion}

La segunda igualdad en la definición de derivada direccional es consecuencia de la
diferenciabilidad de $\mathcal{L}$: la linealización de primer orden de $\mathcal{L}$
alrededor de $\bm{\theta}$ es $\mathcal{L}(\bm{\theta}+h\bm{v}) \approx
\mathcal{L}(\bm{\theta}) + h\langle\nabla_{\bm{\theta}}\mathcal{L}, \bm{v}\rangle$,
de donde se obtiene el resultado dividiendo por $h$ y tomando el límite.

\subsection{El gradiente como dirección de máximo ascenso}
\label{subsec:max_ascenso}

La siguiente proposición es la justificación matemática precisa de todos los
algoritmos de descenso de gradiente que se estudiarán en el Capítulo~5.

\begin{proposicion}{El gradiente apunta en la dirección de máximo ascenso}
\label{prop:max_ascenso}
Sea $\mathcal{L} \in C^1(\mathbb{R}^p)$ y $\bm{\theta} \in \mathbb{R}^p$
con $\nabla_{\bm{\theta}}\mathcal{L}(\bm{\theta}) \neq \bm{0}$. Entonces:
\[
  \max_{\|\bm{v}\|_2=1} D_{\bm{v}}\mathcal{L}(\bm{\theta})
  = \|\nabla_{\bm{\theta}}\mathcal{L}(\bm{\theta})\|_2,
\]
y el máximo se alcanza en la dirección
$\bm{v}^* = \nabla_{\bm{\theta}}\mathcal{L}(\bm{\theta})\,/\,\|\nabla_{\bm{\theta}}\mathcal{L}(\bm{\theta})\|_2$.
En consecuencia, la dirección de \textbf{máximo descenso} es
$-\nabla_{\bm{\theta}}\mathcal{L}(\bm{\theta})\,/\,\|\nabla_{\bm{\theta}}\mathcal{L}(\bm{\theta})\|_2$.
\end{proposicion}

\begin{proof}
Por definición, $D_{\bm{v}}\mathcal{L}(\bm{\theta}) =
\langle\nabla_{\bm{\theta}}\mathcal{L}(\bm{\theta}), \bm{v}\rangle$.
Por la desigualdad de Cauchy-Schwarz (Definición~\ref{def:inner}):
\[
  \langle\nabla_{\bm{\theta}}\mathcal{L}(\bm{\theta}), \bm{v}\rangle
  \;\leq\; \|\nabla_{\bm{\theta}}\mathcal{L}(\bm{\theta})\|_2\,\|\bm{v}\|_2
  = \|\nabla_{\bm{\theta}}\mathcal{L}(\bm{\theta})\|_2,
\]
con igualdad si y solo si $\bm{v}$ es paralelo a
$\nabla_{\bm{\theta}}\mathcal{L}(\bm{\theta})$. Como $\|\bm{v}\|_2 = 1$,
el máximo se alcanza en
$\bm{v}^* = \nabla_{\bm{\theta}}\mathcal{L}(\bm{\theta})\,/\,
\|\nabla_{\bm{\theta}}\mathcal{L}(\bm{\theta})\|_2$.
\end{proof}

\begin{fisicaconexion}{El gradiente en campos escalares físicos}
La Proposición~\ref{prop:max_ascenso} es un resultado que el físico ya conoce
en otro lenguaje. El gradiente de un campo escalar apunta en la dirección de
máximo crecimiento: el gradiente de la temperatura $\nabla T(\bm{r})$ apunta
hacia el lugar más caliente; el gradiente del potencial eléctrico $\nabla V(\bm{r})$
apunta hacia el lugar de mayor potencial; el gradiente de la presión $\nabla P(\bm{r})$
apunta hacia la zona de mayor presión. En todos estos casos, la fuerza física
asociada es proporcional al \emph{negativo} del gradiente: la corriente de calor
fluye de mayor a menor temperatura, el campo eléctrico $\bm{E} = -\nabla V$ apunta
de mayor a menor potencial, el fluido se mueve hacia menor presión. El descenso de
gradiente en optimización es exactamente esa misma idea: moverse en la dirección
$-\nabla\mathcal{L}$ como si la función de pérdida fuera un potencial que genera
una ``fuerza'' que empuja al sistema hacia el mínimo.
\end{fisicaconexion}

\begin{observacion}{Localidad del gradiente}
La Proposición~\ref{prop:max_ascenso} afirma que $-\nabla_{\bm{\theta}}\mathcal{L}$
es la mejor dirección de descenso \emph{localmente}: en el punto $\bm{\theta}$,
para desplazamientos infinitesimales. Esto no implica que moverse en esa dirección
indefinidamente sea óptimo. El gradiente cambia a lo largo del camino de descenso,
y la dirección óptima a cada instante puede diferir significativamente de la
dirección inicial. Esta distinción es la fuente de los distintos comportamientos
de los algoritmos de descenso según la tasa de aprendizaje elegida, como se
analizará en el Capítulo~5.
\end{observacion}

\subsection{El jacobiano y el gradiente de la pérdida empírica}
\label{subsec:jacobiano}

Para el cálculo explícito del gradiente de $\mathcal{L}$ es útil introducir el
jacobiano de la familia respecto a los parámetros.

\begin{definicion}{Jacobiano de la familia paramétrica}
\label{def:jacobiano}
Sea $f:\mathbb{R}^n\times\mathbb{R}^p\to\mathbb{R}$ una familia paramétrica de
clase $C^1$ en $\bm{\theta}$. El \textbf{jacobiano} de $f$ evaluado en el conjunto
de datos $\mathcal{D}$ es la matriz $\bm{J}_f(\bm{\theta}) \in \mathbb{R}^{N\times p}$
con entradas:
\[
  \bigl(\bm{J}_f(\bm{\theta})\bigr)_{ij}
  = \frac{\partial f(\bm{x}_i;\bm{\theta})}{\partial\theta_j},
  \qquad i=1,\ldots,N,\quad j=1,\ldots,p.
\]
La $i$-ésima fila de $\bm{J}_f$ es el gradiente de $f(\bm{x}_i;\cdot)$ respecto
a $\bm{\theta}$, traspuesto.
\end{definicion}

\begin{proposicion}{Gradiente de la pérdida MSE}
\label{prop:grad_mse}
Para la pérdida $\ell_{\mathrm{MSE}}$, el gradiente de la pérdida empírica es:
\[
  \nabla_{\bm{\theta}}\mathcal{L}(\bm{\theta})
  = \frac{2}{N}\,\bm{J}_f(\bm{\theta})^\top\,\bigl(\bm{f}(\bm{\theta}) - \bm{y}\bigr),
\]
donde $\bm{f}(\bm{\theta}) = (f(\bm{x}_1;\bm{\theta}),\ldots,f(\bm{x}_N;\bm{\theta}))^\top
\in \mathbb{R}^N$ es el vector de predicciones y $\bm{y} = (y_1,\ldots,y_N)^\top$.
\end{proposicion}

\begin{proof}
Aplicando la regla de la cadena a cada término de la suma:
\[
  \frac{\partial}{\partial\theta_j}(f(\bm{x}_i;\bm{\theta})-y_i)^2
  = 2\,(f(\bm{x}_i;\bm{\theta})-y_i)\,\frac{\partial f(\bm{x}_i;\bm{\theta})}{\partial\theta_j}.
\]
Sumando sobre $i$ y dividiendo por $N$:
\[
  \frac{\partial\mathcal{L}}{\partial\theta_j}
  = \frac{2}{N}\sum_{i=1}^N (f(\bm{x}_i;\bm{\theta})-y_i)
    \,\bigl(\bm{J}_f(\bm{\theta})\bigr)_{ij}.
\]
Reconociendo la estructura de producto matricial,
$\nabla_{\bm{\theta}}\mathcal{L} = \frac{2}{N}\bm{J}_f^\top(\bm{f}-\bm{y})$.
\end{proof}

\begin{observacion}{Recuperación de las ecuaciones normales}
Para la familia lineal, $f(\bm{x}_i;\bm{\theta}) = \bm{v}_i^\top\bm{\theta}$ donde
$\bm{v}_i^\top$ es la $i$-ésima fila de $\bm{V}$. Entonces $\bm{J}_f(\bm{\theta})
= \bm{V}$ (el jacobiano es la matriz de diseño, constante en $\bm{\theta}$) y la
Proposición~\ref{prop:grad_mse} da:
\[
  \nabla_{\bm{\theta}}\mathcal{L}(\bm{\theta})
  = \frac{2}{N}\bm{V}^\top(\bm{V}\bm{\theta}-\bm{y}).
\]
Igualando a cero se obtienen exactamente las ecuaciones normales de la
Proposición~\ref{prop:normales}. Para una familia no lineal, $\bm{J}_f$ depende
de $\bm{\theta}$ y el sistema $\nabla_{\bm{\theta}}\mathcal{L} = \bm{0}$ es
no lineal: no tiene solución de forma cerrada en general.
\end{observacion}

%───────────────────────────────────────────────────────────────────────────────
\section{El hessiano: curvatura local del paisaje}
\label{sec:hessiano}
%───────────────────────────────────────────────────────────────────────────────

\dificultad{3}{3}

El gradiente describe el cambio de primer orden de $\mathcal{L}$. El hessiano
describe el cambio de segundo orden: la curvatura. Esta información es la que
distingue entre un mínimo, un máximo y un punto de silla.

\begin{definicion}{Matriz hessiana}
\label{def:hessiano}
Sea $\mathcal{L} \in C^2(\mathbb{R}^p)$. La \textbf{matriz hessiana} de
$\mathcal{L}$ en $\bm{\theta}$ es la matriz $\bm{H}(\bm{\theta}) \in
\mathbb{R}^{p\times p}$ con entradas:
\[
  H_{ij}(\bm{\theta})
  = \frac{\partial^2\mathcal{L}}{\partial\theta_i\,\partial\theta_j}(\bm{\theta}),
  \qquad i,j = 1,\ldots,p.
\]
Por el teorema de Schwarz, $H_{ij} = H_{ji}$ siempre que $\mathcal{L} \in C^2$,
de modo que $\bm{H}(\bm{\theta})$ es simétrica para todo $\bm{\theta}$.
\end{definicion}

La importancia del hessiano proviene de la aproximación de Taylor de segundo orden.

\begin{proposicion}{Aproximación cuadrática local}
\label{prop:taylor2}
Sea $\mathcal{L} \in C^2(\mathbb{R}^p)$, $\bm{\theta}_0 \in \mathbb{R}^p$ y
$\bm{d} \in \mathbb{R}^p$. Entonces:
\[
  \mathcal{L}(\bm{\theta}_0 + \bm{d})
  = \mathcal{L}(\bm{\theta}_0)
  + \langle\nabla_{\bm{\theta}}\mathcal{L}(\bm{\theta}_0),\,\bm{d}\rangle
  + \frac{1}{2}\,\bm{d}^\top\bm{H}(\bm{\theta}_0)\bm{d}
  + O(\|\bm{d}\|_2^3).
\]
\end{proposicion}

Esta expansión es el punto de partida de todos los métodos de optimización de
segundo orden, en particular del método de Newton que se estudiará en el
Capítulo~5.

\subsection{Curvatura direccional}
\label{subsec:curvatura}

\dificultad{3}{3}

En la dirección de un vector unitario $\bm{v}$, la curvatura de $\mathcal{L}$
en $\bm{\theta}_0$ es:
\[
  \kappa_{\bm{v}}(\bm{\theta}_0)
  = \bm{v}^\top\bm{H}(\bm{\theta}_0)\bm{v}.
\]
Cuando $\bm{v} = \bm{e}_j$ (el $j$-ésimo vector canónico), esto es simplemente
$H_{jj}(\bm{\theta}_0)$, la curvatura a lo largo del eje $\theta_j$.

Como $\bm{H}(\bm{\theta}_0)$ es simétrica, el teorema espectral garantiza que
existe una base ortonormal de $\mathbb{R}^p$ formada por autovectores
$\bm{v}_1,\ldots,\bm{v}_p$ con autovalores reales $\mu_1 \geq \mu_2 \geq \cdots
\geq \mu_p$. En estas direcciones la aproximación cuadrática se desacopla:
\[
  \mathcal{L}(\bm{\theta}_0 + \bm{d})
  \approx \mathcal{L}(\bm{\theta}_0)
  + \sum_{k=1}^p \Bigl(
    \langle\nabla_{\bm{\theta}}\mathcal{L}(\bm{\theta}_0),\bm{v}_k\rangle\,d_k
    + \frac{\mu_k}{2}\,d_k^2\Bigr),
\]
donde $d_k = \langle\bm{d},\bm{v}_k\rangle$. La función se comporta localmente
como una suma de parábolas independientes en cada dirección espectral, con curvatura
$\mu_k$: positiva para $\mu_k > 0$ (cóncava hacia arriba, mínimo local en esa
dirección), negativa para $\mu_k < 0$ (cóncava hacia abajo, máximo local), y nula
para $\mu_k = 0$ (plana en esa dirección).

\begin{fisicaconexion}{El hessiano como matriz de modos normales}
En mecánica clásica, el hessiano del potencial de energía $V(\bm{q})$ evaluado
en un punto de equilibrio $\bm{q}_0$ (donde $\nabla V(\bm{q}_0) = \bm{0}$) es
exactamente la \textbf{matriz de constantes de fuerza}: la aproximación armónica
$V(\bm{q}_0+\bm{u}) \approx V(\bm{q}_0) + \frac{1}{2}\bm{u}^\top\bm{H}_V\bm{u}$
define un sistema de osciladores acoplados. Diagonalizar $\bm{H}_V$ es precisamente
encontrar los \textbf{modos normales}: los autovectores son las direcciones de
vibración desacopladas y los autovalores son las constantes de rigidez en cada
dirección, proporcionales al cuadrado de las frecuencias propias
($\omega_k^2 \propto \mu_k$). Un modo con $\mu_k > 0$ corresponde a una dirección
estable (oscilación); uno con $\mu_k < 0$ corresponde a una dirección inestable
(la partícula se aleja del equilibrio, análogo a un punto de silla). El análisis
del hessiano de la función de pérdida en un punto crítico es, formalmente, la misma
operación que el análisis de modos normales en mecánica.
\end{fisicaconexion}

\subsection{Hessiano de la pérdida MSE}
\label{subsec:hess_mse}

\begin{proposicion}{Hessiano de la pérdida MSE}
\label{prop:hess_mse}
Para la pérdida $\ell_{\mathrm{MSE}}$, el hessiano de la pérdida empírica es:
\[
  \bm{H}(\bm{\theta})
  = \frac{2}{N}\Bigl[\bm{J}_f(\bm{\theta})^\top\bm{J}_f(\bm{\theta})
    + \sum_{i=1}^N r_i(\bm{\theta})\,\bm{G}_i(\bm{\theta})\Bigr],
\]
donde $r_i(\bm{\theta}) = f(\bm{x}_i;\bm{\theta}) - y_i$ son los \textbf{residuos}
y $\bm{G}_i(\bm{\theta}) \in \mathbb{R}^{p\times p}$ es la matriz hessiana de
$f(\bm{x}_i;\cdot)$ en $\bm{\theta}$.
\end{proposicion}

\begin{proof}
De la Proposición~\ref{prop:grad_mse}, el $j$-ésimo componente del gradiente es:
\[
  \frac{\partial\mathcal{L}}{\partial\theta_j}
  = \frac{2}{N}\sum_{i=1}^N r_i\,\frac{\partial f(\bm{x}_i;\bm{\theta})}{\partial\theta_j}.
\]
Derivando respecto a $\theta_k$:
\begin{align*}
  H_{jk} = \frac{\partial^2\mathcal{L}}{\partial\theta_k\partial\theta_j}
  &= \frac{2}{N}\sum_{i=1}^N \left[
    \frac{\partial f(\bm{x}_i;\bm{\theta})}{\partial\theta_k}
    \frac{\partial f(\bm{x}_i;\bm{\theta})}{\partial\theta_j}
    + r_i\,\frac{\partial^2 f(\bm{x}_i;\bm{\theta})}{\partial\theta_k\partial\theta_j}
  \right].
\end{align*}
Reconociendo $(\bm{J}_f^\top\bm{J}_f)_{jk} = \sum_i
\frac{\partial f_i}{\partial\theta_j}\frac{\partial f_i}{\partial\theta_k}$
y $(\bm{G}_i)_{jk} = \frac{\partial^2 f_i}{\partial\theta_j\partial\theta_k}$,
se obtiene el resultado.
\end{proof}

\begin{observacion}{La aproximación de Gauss-Newton}
\label{obs:gaussnewton}
En el mínimo $\bm{\theta}^*$, si el ajuste es bueno, los residuos $r_i(\bm{\theta}^*)
\approx 0$ y el segundo término del hessiano es despreciable:
\[
  \bm{H}(\bm{\theta}^*) \approx \frac{2}{N}\bm{J}_f(\bm{\theta}^*)^\top\bm{J}_f(\bm{\theta}^*).
\]
Esta aproximación, conocida como \textbf{aproximación de Gauss-Newton}, es semidefinida
positiva (pues $\bm{J}_f^\top\bm{J}_f$ siempre lo es) y es de fácil cómputo porque
no requiere calcular las segundas derivadas de $f$. Para la familia lineal,
$\bm{J}_f = \bm{V}$ y el hessiano es exactamente $\frac{2}{N}\bm{V}^\top\bm{V}$,
constante e independiente de $\bm{\theta}$: el paisaje cuadrático del Capítulo~2
emerge como caso particular.
\end{observacion}

%───────────────────────────────────────────────────────────────────────────────
\section{Puntos críticos y su clasificación}
\label{sec:criticos}
%───────────────────────────────────────────────────────────────────────────────

\dificultad{2}{3}

\begin{definicion}{Punto crítico}
\label{def:critico}
Un punto $\bm{\theta}^* \in \mathbb{R}^p$ es un \textbf{punto crítico} de
$\mathcal{L} \in C^1(\mathbb{R}^p)$ si:
\[
  \nabla_{\bm{\theta}}\mathcal{L}(\bm{\theta}^*) = \bm{0}.
\]
\end{definicion}

Los puntos críticos son los candidatos a mínimos de $\mathcal{L}$, pero no todos
los puntos críticos son mínimos.

\begin{proposicion}{Condición necesaria de primer orden}
\label{prop:cond1}
Si $\bm{\theta}^*$ es un mínimo local de $\mathcal{L} \in C^1(\mathbb{R}^p)$,
entonces $\bm{\theta}^*$ es un punto crítico.
\end{proposicion}

\begin{proof}
Supóngase que $\bm{\theta}^*$ es mínimo local pero $\nabla_{\bm{\theta}}\mathcal{L}(\bm{\theta}^*)
\neq \bm{0}$. Sea $\bm{v} = -\nabla_{\bm{\theta}}\mathcal{L}(\bm{\theta}^*)/
\|\nabla_{\bm{\theta}}\mathcal{L}(\bm{\theta}^*)\|_2$. Por diferenciabilidad:
\[
  \mathcal{L}(\bm{\theta}^* + h\bm{v})
  = \mathcal{L}(\bm{\theta}^*)
  + h\underbrace{\langle\nabla_{\bm{\theta}}\mathcal{L}(\bm{\theta}^*),\bm{v}\rangle}_{-\|\nabla\mathcal{L}\|_2 < 0}
  + O(h^2).
\]
Para $h > 0$ suficientemente pequeño, el lado derecho es menor que
$\mathcal{L}(\bm{\theta}^*)$, contradiciendo que $\bm{\theta}^*$ es mínimo local.
\end{proof}

\subsection{Clasificación mediante el hessiano}
\label{subsec:clasificacion}

\begin{proposicion}{Condición suficiente de segundo orden}
\label{prop:cond2}
Sea $\bm{\theta}^*$ un punto crítico de $\mathcal{L} \in C^2(\mathbb{R}^p)$.
\begin{enumerate}
  \item Si $\bm{H}(\bm{\theta}^*)$ es \textbf{definida positiva} ($\mu_k > 0$
        para todo $k$), entonces $\bm{\theta}^*$ es un \textbf{mínimo local
        estricto}.
  \item Si $\bm{H}(\bm{\theta}^*)$ es \textbf{definida negativa} ($\mu_k < 0$
        para todo $k$), entonces $\bm{\theta}^*$ es un \textbf{máximo local
        estricto}.
  \item Si $\bm{H}(\bm{\theta}^*)$ es \textbf{indefinida} (tiene autovalores
        positivos y negativos), entonces $\bm{\theta}^*$ es un \textbf{punto
        de silla}.
\end{enumerate}
\end{proposicion}

\begin{proof}
Para el caso 1: por la Proposición~\ref{prop:taylor2}, con $\nabla\mathcal{L}(\bm{\theta}^*)
= \bm{0}$:
\[
  \mathcal{L}(\bm{\theta}^* + \bm{d})
  = \mathcal{L}(\bm{\theta}^*)
  + \frac{1}{2}\bm{d}^\top\bm{H}(\bm{\theta}^*)\bm{d}
  + O(\|\bm{d}\|_2^3).
\]
Si $\bm{H}(\bm{\theta}^*)$ es definida positiva con mínimo autovalor $\mu_{\min} > 0$:
\[
  \bm{d}^\top\bm{H}(\bm{\theta}^*)\bm{d} \geq \mu_{\min}\|\bm{d}\|_2^2.
\]
Para $\|\bm{d}\|_2$ suficientemente pequeño, el término $O(\|\bm{d}\|_2^3)$ es
dominado por $\frac{\mu_{\min}}{2}\|\bm{d}\|_2^2 > 0$, por tanto
$\mathcal{L}(\bm{\theta}^*+\bm{d}) > \mathcal{L}(\bm{\theta}^*)$.
Los casos 2 y 3 se demuestran de forma análoga.
\end{proof}

\begin{definicion}{Índice de Morse de un punto crítico}
\label{def:morse}
El \textbf{índice de Morse} de un punto crítico no degenerado $\bm{\theta}^*$
(donde $\bm{H}(\bm{\theta}^*)$ no tiene autovalores nulos) es el número de
autovalores negativos de $\bm{H}(\bm{\theta}^*)$:
\[
  \mathrm{ind}(\bm{\theta}^*) = \#\{k : \mu_k(\bm{H}(\bm{\theta}^*)) < 0\}.
\]
Un mínimo local tiene índice $0$. Un máximo local tiene índice $p$. Un punto de
silla de índice $q \in \{1,\ldots,p-1\}$ tiene $q$ direcciones de descenso y
$p-q$ de ascenso.
\end{definicion}

\subsection{Puntos de silla en alta dimensión}
\label{subsec:sillas_alta_dim}

\dificultad{3}{3}

En $\mathbb{R}^1$, los puntos críticos son mínimos, máximos o puntos de inflexión.
En $\mathbb{R}^p$ con $p$ grande, la situación es radicalmente diferente.

\begin{observacion}{Abundancia de puntos de silla}
\label{obs:sillas}
Un punto crítico aleatorio de una función ``genérica'' en $\mathbb{R}^p$ tiene cada
autovalor del hessiano positivo o negativo con probabilidad $1/2$,
independientemente. La probabilidad de que sea un mínimo (todos los $p$ autovalores
positivos) es $2^{-p}$. Para $p = 10^6$ esta probabilidad es $\approx 10^{-300000}$:
en la práctica, los únicos puntos críticos que existen son puntos de silla.

Esto tiene consecuencias importantes para la optimización en redes neuronales.
Contrariamente a la creencia inicial, el problema principal no es quedar atrapado
en mínimos locales: los mínimos locales son exponencialmente raros. El verdadero
obstáculo son los puntos de silla, donde el gradiente es cero pero el punto no
es un mínimo. Los algoritmos de descenso de gradiente con ruido (como SGD) tienen
la propiedad de escapar de los puntos de silla porque las fluctuaciones estocásticas
rompen la simetría local. Esto se analizará en el Capítulo~5.
\end{observacion}

\begin{fisicaconexion}{Puntos de silla en superficies de energía potencial}
Los puntos de silla son objetos centrales en física y química desde antes del
aprendizaje automático. En física de sólidos y química cuántica, el \textbf{estado
de transición} de una reacción química es un punto de silla sobre la
\textbf{superficie de energía potencial} (PES): es un punto de mínima energía en
todas las direcciones salvo una (la coordenada de reacción), en la cual es un
máximo. El índice de Morse es exactamente $1$. La tasa de reacción depende
exponencialmente de la altura de esa barrera, según la teoría del estado de
transición (Eyring, 1935). En espectroscopía de moléculas poliatómicas, los
puntos de silla de la PES determinan los canales de disociación y la estabilidad
de los estados vibracionales. El mismo problema matemático —encontrar y clasificar
puntos críticos de una función escalar en alta dimensión— aparece tanto en la
optimización de redes neuronales como en el estudio de la reactividad química.
\end{fisicaconexion}

%───────────────────────────────────────────────────────────────────────────────
\section{Convexidad: el caso ideal}
\label{sec:convexidad}
%───────────────────────────────────────────────────────────────────────────────

\dificultad{2}{3}

\begin{definicion}{Función convexa}
\label{def:convexa}
Una función $\mathcal{L}:\mathbb{R}^p \to \mathbb{R}$ es \textbf{convexa} si para
todo $\bm{\theta}_1, \bm{\theta}_2 \in \mathbb{R}^p$ y todo $t \in [0,1]$:
\[
  \mathcal{L}\bigl(t\bm{\theta}_1 + (1-t)\bm{\theta}_2\bigr)
  \;\leq\; t\,\mathcal{L}(\bm{\theta}_1) + (1-t)\,\mathcal{L}(\bm{\theta}_2).
\]
Es \textbf{estrictamente convexa} si la desigualdad es estricta para $\bm{\theta}_1
\neq \bm{\theta}_2$ y $t \in (0,1)$.
\end{definicion}

La caracterización diferencial de la convexidad conecta directamente con el hessiano.

\begin{proposicion}{Caracterización diferencial de la convexidad}
\label{prop:conv_diff}
Sea $\mathcal{L} \in C^2(\mathbb{R}^p)$. Entonces:
\begin{enumerate}
  \item $\mathcal{L}$ es convexa si y solo si $\bm{H}(\bm{\theta}) \succeq 0$
        para todo $\bm{\theta} \in \mathbb{R}^p$ (hessiano semidefinido positivo
        en todo punto).
  \item $\mathcal{L}$ es estrictamente convexa si $\bm{H}(\bm{\theta}) \succ 0$
        para todo $\bm{\theta}$ (hessiano definido positivo en todo punto).
\end{enumerate}
\end{proposicion}

La demostración usa el teorema fundamental del cálculo para expresar la diferencia
$\mathcal{L}(t\bm{\theta}_1+(1-t)\bm{\theta}_2) - t\mathcal{L}(\bm{\theta}_1)
- (1-t)\mathcal{L}(\bm{\theta}_2)$ como una integral que involucra el hessiano;
se deja como Ejercicio Resuelto~\ref{ej:conv_hess}.

\begin{proposicion}{La pérdida MSE lineal es convexa}
\label{prop:mse_convexa}
Para la familia lineal en parámetros con pérdida MSE,
$\mathcal{L}(\bm{\theta}) = \frac{1}{N}\|\bm{V}\bm{\theta}-\bm{y}\|_2^2$ es
convexa. Es estrictamente convexa si y solo si $\rank(\bm{V}) = p$.
\end{proposicion}

\begin{proof}
El hessiano es $\bm{H} = \frac{2}{N}\bm{V}^\top\bm{V}$, constante. Para todo
$\bm{d} \in \mathbb{R}^p$:
\[
  \bm{d}^\top\bm{H}\bm{d}
  = \frac{2}{N}\bm{d}^\top\bm{V}^\top\bm{V}\bm{d}
  = \frac{2}{N}\|\bm{V}\bm{d}\|_2^2 \geq 0,
\]
con igualdad si y solo si $\bm{V}\bm{d} = \bm{0}$, es decir, $\bm{d} \in
\ker(\bm{V})$. Si $\rank(\bm{V}) = p$, el núcleo es trivial ($\bm{d} = \bm{0}$)
y $\bm{H}$ es definida positiva: la función es estrictamente convexa.
\end{proof}

\begin{proposicion}{La regularización restaura la convexidad estricta}
\label{prop:reg_convexa}
Para todo $\lambda > 0$, la función $F_\lambda(\bm{\theta}) = \mathcal{L}(\bm{\theta})
+ \lambda\|\bm{\theta}\|_2^2$ es estrictamente convexa, independientemente del
rango de $\bm{V}$.
\end{proposicion}

\begin{proof}
El hessiano de $F_\lambda$ es $\bm{H} + 2\lambda\bm{I}$. Para todo $\bm{d}
\neq \bm{0}$:
\[
  \bm{d}^\top(\bm{H}+2\lambda\bm{I})\bm{d}
  = \bm{d}^\top\bm{H}\bm{d} + 2\lambda\|\bm{d}\|_2^2
  \geq 2\lambda\|\bm{d}\|_2^2 > 0.
\]
\end{proof}

\subsection{Consecuencias de la convexidad para la optimización}
\label{subsec:conv_optim}

La convexidad es una propiedad extraordinariamente útil en optimización porque
elimina la posibilidad de mínimos locales que no sean globales.

\begin{proposicion}{Todo punto crítico de una función convexa es mínimo global}
\label{prop:critico_global}
Sea $\mathcal{L}:\mathbb{R}^p\to\mathbb{R}$ convexa y diferenciable.
Si $\nabla_{\bm{\theta}}\mathcal{L}(\bm{\theta}^*) = \bm{0}$, entonces
$\bm{\theta}^*$ es un mínimo global.
\end{proposicion}

\begin{proof}
Por convexidad y diferenciabilidad, para todo $\bm{\theta} \in \mathbb{R}^p$:
\[
  \mathcal{L}(\bm{\theta})
  \geq \mathcal{L}(\bm{\theta}^*)
    + \langle\nabla_{\bm{\theta}}\mathcal{L}(\bm{\theta}^*),\,\bm{\theta}-\bm{\theta}^*\rangle
  = \mathcal{L}(\bm{\theta}^*),
\]
donde la primera desigualdad es la condición de primer orden de convexidad
(la función yace por encima de cualquier hiperplano tangente) y la igualdad
usa $\nabla\mathcal{L}(\bm{\theta}^*) = \bm{0}$.
\end{proof}

Para una función estrictamente convexa, el mínimo global es único: si existieran
dos mínimos $\bm{\theta}_1^* \neq \bm{\theta}_2^*$, la convexidad estricta daría
$\mathcal{L}(\frac{1}{2}(\bm{\theta}_1^*+\bm{\theta}_2^*)) <
\frac{1}{2}\mathcal{L}(\bm{\theta}_1^*) + \frac{1}{2}\mathcal{L}(\bm{\theta}_2^*)
= \mathcal{L}(\bm{\theta}_1^*)$, contradiciendo la minimalidad.

\begin{fisicaconexion}{Convexidad y equilibrio estable en mecánica}
La convexidad del potencial de energía es la condición matemática de
\textbf{equilibrio estable}. Un potencial $V(\bm{q})$ convexo implica que todo
punto de equilibrio (donde $\nabla V = \bm{0}$) es un mínimo global: si se
perturba el sistema, la fuerza restauradora $-\nabla V$ lo devuelve al mismo
estado de equilibrio, sin importar cuán grande sea la perturbación. El potencial
armónico $V = \frac{1}{2}k\|\bm{q}\|^2$ es el ejemplo arquetípico: estrictamente
convexo, con un único mínimo global en $\bm{q} = \bm{0}$, y toda trayectoria
dinámica converge a él (con disipación). La pérdida MSE lineal tiene exactamente
esta geometría: es el análogo computacional de un oscilador armónico en el espacio
de parámetros. Cuando el potencial deja de ser convexo ---un doble pozo, por
ejemplo--- aparecen múltiples equilibrios y el estado final depende de las
condiciones iniciales. Las redes neuronales están en ese régimen.
\end{fisicaconexion}

%───────────────────────────────────────────────────────────────────────────────
\section{No convexidad: el paisaje real de las redes neuronales}
\label{sec:noconvexidad}
%───────────────────────────────────────────────────────────────────────────────

\dificultad{2}{2}

La convexidad es la excepción, no la regla. En cuanto la familia paramétrica es
no lineal, la pérdida empírica pierde convexidad. Para las redes neuronales, esto
no es un accidente numérico: la no convexidad está inscrita en la estructura
arquitectónica del modelo.

\subsection{La pérdida de una red neuronal es no convexa}
\label{subsec:noconv_red}

Se muestra la no convexidad con el modelo más simple posible.

\begin{ejemplo}{No convexidad de una red de una neurona}
\label{ex:noconvex}
Considérese la familia de una sola neurona con activación sigmoide $\sigma(z) =
(1+e^{-z})^{-1}$:
\[
  f(x; w, b) = \sigma(wx + b), \qquad \bm{\theta} = (w,b)^\top \in \mathbb{R}^2.
\]
Con el dato único $(x_1, y_1) = (0, 0.7)$ y pérdida MSE:
\[
  \mathcal{L}(w,b) = \bigl(\sigma(b) - 0.7\bigr)^2.
\]
Esta función no depende de $w$: es constante en la dirección $w$, y por tanto
tiene $\partial^2\mathcal{L}/\partial w^2 = 0$. El hessiano en todo punto tiene
un autovalor cero (dirección plana) y no es definido positivo. Cualquier desplazamiento
en la dirección $w$ no modifica $\mathcal{L}$: el mínimo no es un punto aislado
sino una recta entera $\{(w, b^*) : w \in \mathbb{R}\}$ donde $b^* = \sigma^{-1}(0.7)$.
\end{ejemplo}

El ejemplo anterior muestra que incluso la red más simple tiene simetrías que
producen degeneración. Para redes con más neuronas, las simetrías son más ricas.

\subsection{Simetrías y mínimos equivalentes}
\label{subsec:simetrias}

\begin{observacion}{Simetría de permutación}
\label{obs:perm}
Sea una red neuronal con una capa oculta de $m$ neuronas:
\[
  f(\bm{x};\bm{\theta}) = \sum_{j=1}^m v_j\,\sigma(\bm{w}_j^\top\bm{x} + b_j).
\]
Para cualquier permutación $\pi$ de $\{1,\ldots,m\}$, el modelo con parámetros
permutados $(v_{\pi(j)}, \bm{w}_{\pi(j)}, b_{\pi(j)})$ produce exactamente la
misma función. Por tanto, si $\bm{\theta}^*$ es un mínimo de $\mathcal{L}$,
también lo son sus $m!$ permutaciones. Una red con $m = 100$ neuronas tiene al
menos $100! \approx 10^{158}$ mínimos equivalentes.
\end{observacion}

\begin{observacion}{Simetría de escala}
\label{obs:escala}
Para activaciones homogéneas de grado $d$ (como ReLU, que satisface
$\sigma(\alpha z) = \alpha\sigma(z)$ para $\alpha > 0$), existe la simetría de
escala: reemplazar $\bm{w}_j \to \alpha\bm{w}_j$ y $v_j \to v_j/\alpha$ deja
invariante la función $f$. Esto crea valles planos en el paisaje de $\mathcal{L}$:
curvas de parámetros equivalentes que no son mínimos aislados.
\end{observacion}

Estas simetrías implican que el paisaje de pérdida de una red neuronal está lejos
de ser un conjunto discreto de puntos críticos aislados. La estructura real es mucho
más rica: manifolds de parámetros equivalentes, valles elongados y regiones planas
que los algoritmos de gradiente recorren sin descender.

\subsection{El fenómeno de \textit{double descent} y los mínimos locales en alta dimensión}
\label{subsec:double_descent}

\dificultad{3}{2}

Un resultado sorprendente, establecido empíricamente y justificado parcialmente
de forma teórica en los últimos años, sugiere que la imagen clásica de ``muchos
mínimos locales malos'' no describe bien el comportamiento real de las redes
sobredeterminadas.

\begin{observacion}{Calidad de los mínimos locales en redes sobredeterminadas}
\label{obs:minimos_locales}
Para redes neuronales con más parámetros que datos ($p \gg N$), los resultados
empíricos y teóricos parciales indican que los mínimos locales de $\mathcal{L}$
tienden a tener valores similares al mínimo global. Esta propiedad, que no tiene
análogo en funciones arbitrarias, parece estar relacionada con la sobredeterminación
del sistema: cuando hay muchas más variables que restricciones, la pérdida es
aproximadamente plana en las direcciones adicionales y los mínimos son
``equivalentemente buenos.''

Relacionado con esto, el \textbf{fenómeno de \textit{double descent}}
\cite{belkin2019reconciling} describe el comportamiento del error de generalización
como función de la capacidad $p$: decrece hasta un mínimo (como predice la curva
clásica del Capítulo~2), aumenta en la zona de interpolación exacta ($p \approx N$),
y luego \emph{vuelve a decrecer} al entrar en el régimen sobredeterminado ($p \gg N$).
Este comportamiento contradice el análisis sesgo-varianza clásico del Capítulo~2
y sugiere que el régimen sobredeterminado tiene propiedades de regularización
implícitas que aún no se comprenden completamente.

Este fenómeno se analizará en profundidad en el Capítulo~8, una vez que tengamos
la arquitectura completa y las herramientas de análisis de capacidad necesarias
para formularlo con precisión. Lo que importa retener aquí es que la no convexidad
del paisaje de las redes neuronales no implica necesariamente dificultades
prácticas en la optimización: en el régimen sobredeterminado, la estructura del
paisaje es más benigna de lo que el análisis clásico de optimización sugiere.
\end{observacion}

%───────────────────────────────────────────────────────────────────────────────
\section{Condicionamiento y velocidad de convergencia}
\label{sec:condicionamiento}
%───────────────────────────────────────────────────────────────────────────────

\dificultad{3}{3}

El espectro del hessiano no solo determina si un punto crítico es mínimo o silla:
también cuantifica la dificultad de la optimización en la vecindad de ese punto.

\begin{definicion}{Número de condicionamiento del hessiano}
\label{def:cond_hess}
Sea $\bm{H}(\bm{\theta})$ un hessiano definido positivo con autovalores $\mu_1
\geq \mu_2 \geq \cdots \geq \mu_p > 0$. El \textbf{número de condicionamiento}
del hessiano en $\bm{\theta}$ es:
\[
  \kappa(\bm{H}(\bm{\theta})) = \frac{\mu_1}{\mu_p}.
\]
\end{definicion}

La relación entre el condicionamiento del hessiano y el condicionamiento de la
matriz de diseño del Capítulo~2 es directa: para la familia lineal,
$\bm{H} = \frac{2}{N}\bm{V}^\top\bm{V}$ cuyos autovalores son $\frac{2}{N}\sigma_i^2$,
de modo que $\kappa(\bm{H}) = (\sigma_1/\sigma_p)^2 = \kappa(\bm{V})^2$.
Un sistema que era difícil de invertir en el Capítulo~2 es doblemente difícil de
optimizar.

\subsection{Convergencia del descenso de gradiente para funciones cuadráticas}
\label{subsec:conv_gradiente}

Para el caso exactamente cuadrático (familia lineal) se puede analizar la
convergencia de forma precisa. El resultado siguiente anticipa el análisis
del Capítulo~5 y motiva la búsqueda de algoritmos que mitiguen el problema
del condicionamiento.

\begin{proposicion}{Convergencia del gradiente para pérdidas cuadráticas}
\label{prop:conv_cuad}
Sea $\mathcal{L}(\bm{\theta}) = \frac{1}{N}\|\bm{V}\bm{\theta}-\bm{y}\|_2^2$ con
$\bm{V}$ de rango completo, y sea $\bm{H} = \frac{2}{N}\bm{V}^\top\bm{V}$. El
método de descenso de gradiente con tasa de aprendizaje $\eta$:
\[
  \bm{\theta}^{(k+1)} = \bm{\theta}^{(k)} - \eta\,\nabla_{\bm{\theta}}\mathcal{L}(\bm{\theta}^{(k)})
\]
converge al mínimo $\bm{\theta}^*$ si y solo si $0 < \eta < 2/\mu_1$. Con la
tasa óptima $\eta^* = 2/(\mu_1+\mu_p)$, el error en cada iteración satisface:
\[
  \|\bm{\theta}^{(k)} - \bm{\theta}^*\|_2
  \;\leq\; \left(\frac{\kappa-1}{\kappa+1}\right)^k \|\bm{\theta}^{(0)}-\bm{\theta}^*\|_2,
\]
donde $\kappa = \kappa(\bm{H})$.
\end{proposicion}

La demostración se desarrolla en el Ejercicio Resuelto~\ref{ej:conv_cuad}. El
resultado tiene una consecuencia cuantitativa inmediata: para reducir el error
en un factor $e^{-1}$, se necesitan $k \approx (\kappa+1)/2$ iteraciones. Para
un hessiano bien condicionado ($\kappa = 10$): $k \approx 6$ iteraciones. Para
un hessiano mal condicionado ($\kappa = 10^4$): $k \approx 5000$ iteraciones.
Este es el costo concreto del mal condicionamiento en términos computacionales.

\begin{observacion}{El problema de los valles elongados}
\label{obs:valles}
Un hessiano con $\kappa \gg 1$ corresponde geométricamente a superficies de
nivel $\mathcal{L}(\bm{\theta}) = c$ que son elipsoides muy elongados: la función
cae rápidamente en unas direcciones (autovectores de $\mu_1$) y lentamente en
otras ($\mu_p$). El descenso de gradiente en este terreno produce trayectorias en
zigzag que oscilan transversalmente al valle mientras avanzan lentamente hacia el
fondo. La solución es usar información de curvatura (el hessiano o aproximaciones
al mismo) para escalar correctamente las direcciones de descenso. Los optimizadores
adaptativos del Capítulo~5 (Adam, RMSProp) son aproximaciones prácticas a esta idea.
\end{observacion}

\begin{fisicaconexion}{Mal condicionamiento en mediciones experimentales}
El número de condicionamiento $\kappa$ es exactamente el mismo objeto que aparece
en la propagación de errores en física experimental. Cuando se invierte un sistema
lineal $\bm{A}\bm{x} = \bm{b}$ para determinar los parámetros $\bm{x}$ a partir
de mediciones $\bm{b}$, un error $\delta\bm{b}$ en las mediciones se amplifica
en el resultado hasta en un factor $\kappa(\bm{A})$: $\|\delta\bm{x}\|/\|\bm{x}\|
\leq \kappa(\bm{A})\,\|\delta\bm{b}\|/\|\bm{b}\|$. En espectroscopía de Fourier,
el condicionamiento de la matriz de transformada discreta determina cuánto amplifica
el proceso de inversión el ruido del detector. En tomografía, el condicionamiento
de la matriz de proyección determina cuánto ruido introduce la retro-proyección
filtrada. El \textbf{zigzag del descenso de gradiente} en un paisaje mal condicionado
es el análogo iterativo de ese fenómeno: cada paso amplifica las componentes de
ruido en las direcciones de baja curvatura, produciendo oscilaciones transversales
al valle antes de converger al fondo.
\end{fisicaconexion}

%───────────────────────────────────────────────────────────────────────────────
\section{Resumen del Capítulo 3}
\label{sec:resumen3}
%───────────────────────────────────────────────────────────────────────────────

La geometría del espacio de parámetros tiene dos caras, y este capítulo las ha
estudiado con las mismas herramientas.

En el caso convexo ---que corresponde a familias lineales en parámetros--- el
paisaje de $\mathcal{L}$ es una parábola sin trampas: todo punto crítico es un
mínimo global, el hessiano es constante y sus autovalores determinan completamente
la velocidad de cualquier algoritmo de descenso. El condicionamiento de ese hessiano
es el cuadrado del condicionamiento de la matriz de diseño, que ya apareció en el
Capítulo~2 como obstáculo para la inversión directa. Aquí reaparece como obstáculo
para la optimización iterativa.

En el caso no convexo ---familias no lineales, redes neuronales--- el paisaje es
radicalmente más rico. La pérdida puede tener múltiples puntos críticos, y las
simetrías del modelo garantizan que los mínimos vienen en familias equivalentes.
Sin embargo, la alta dimensión introduce un fenómeno no intuitivo: los mínimos
locales de baja calidad son exponencialmente raros, y los verdaderos obstáculos son
los puntos de silla, de los que es posible escapar con perturbaciones estocásticas.

El capítulo cierra con dos promesas hacia adelante. El Capítulo~4 construirá la
arquitectura que crea ese paisaje no convexo: la neurona, la capa y la red profunda.
El Capítulo~5 diseñará los algoritmos para navegarlo: descenso de gradiente, sus
variantes estocásticas y los optimizadores adaptativos que mitigan el problema del
condicionamiento.

%───────────────────────────────────────────────────────────────────────────────
\section*{Ejercicios resueltos}
\addcontentsline{toc}{section}{Ejercicios resueltos}
%───────────────────────────────────────────────────────────────────────────────

\begin{ejercicio}
\label{ej:conv_hess}
\textbf{(Convexidad y hessiano semidefinido positivo).}

Demostrar que si $\mathcal{L} \in C^2(\mathbb{R}^p)$ y $\bm{H}(\bm{\theta})
\succeq 0$ para todo $\bm{\theta}$, entonces $\mathcal{L}$ es convexa.

\medskip\noindent\textit{Solución.}
Sean $\bm{\theta}_1, \bm{\theta}_2 \in \mathbb{R}^p$ y $t \in [0,1]$. Definamos
$\bm{\theta}(s) = \bm{\theta}_1 + s(\bm{\theta}_2 - \bm{\theta}_1)$ para $s \in [0,1]$,
de modo que $\bm{\theta}(0) = \bm{\theta}_1$ y $\bm{\theta}(1) = \bm{\theta}_2$.
Sea $\bm{d} = \bm{\theta}_2 - \bm{\theta}_1$.

Por el teorema fundamental del cálculo:
\[
  \mathcal{L}(\bm{\theta}_2) - \mathcal{L}(\bm{\theta}_1)
  = \int_0^1 \frac{d}{ds}\mathcal{L}(\bm{\theta}(s))\,ds
  = \int_0^1 \langle\nabla_{\bm{\theta}}\mathcal{L}(\bm{\theta}(s)),\bm{d}\rangle\,ds.
\]

Sea $\bm{\theta}^* = t\bm{\theta}_1+(1-t)\bm{\theta}_2$, $\bm{d}_1 = \bm{\theta}_1 -
\bm{\theta}^*$, $\bm{d}_2 = \bm{\theta}_2 - \bm{\theta}^*$. Usando la fórmula
integral de Taylor de segundo orden en cada punto:
\begin{align*}
  t\mathcal{L}(\bm{\theta}_1) + (1-t)\mathcal{L}(\bm{\theta}_2)
  - \mathcal{L}(\bm{\theta}^*)
  &= t\int_0^1(1-s)\bm{d}_1^\top\bm{H}(\bm{\theta}^*+s\bm{d}_1)\bm{d}_1\,ds \\
  &\quad+ (1-t)\int_0^1(1-s)\bm{d}_2^\top\bm{H}(\bm{\theta}^*+s\bm{d}_2)\bm{d}_2\,ds.
\end{align*}
Como $\bm{H} \succeq 0$, ambos integrandos son no negativos, y la diferencia
es no negativa. Esto establece la desigualdad de convexidad. $\square$
\end{ejercicio}

\begin{ejercicio}
\label{ej:conv_cuad}
\textbf{(Convergencia del descenso de gradiente para pérdida cuadrática).}

Sea $\mathcal{L}(\bm{\theta}) = \frac{1}{2}\bm{\theta}^\top\bm{H}\bm{\theta} -
\bm{b}^\top\bm{\theta} + c$ con $\bm{H}$ simétrica definida positiva con
autovalores $\mu_p \leq \cdots \leq \mu_1$. El mínimo es $\bm{\theta}^* = \bm{H}^{-1}\bm{b}$.
Demostrar la cota de convergencia de la Proposición~\ref{prop:conv_cuad}.

\medskip\noindent\textit{Solución.}
Definamos el error $\bm{e}^{(k)} = \bm{\theta}^{(k)} - \bm{\theta}^*$. Usando que
$\nabla\mathcal{L}(\bm{\theta}) = \bm{H}\bm{\theta} - \bm{b} = \bm{H}(\bm{\theta}-\bm{\theta}^*)$:
\[
  \bm{e}^{(k+1)} = \bm{\theta}^{(k+1)} - \bm{\theta}^*
  = \bm{\theta}^{(k)} - \eta\bm{H}\bm{\theta}^{(k)} + \eta\bm{b} - \bm{\theta}^*
  = (\bm{I} - \eta\bm{H})\bm{e}^{(k)}.
\]
Iterando: $\bm{e}^{(k)} = (\bm{I}-\eta\bm{H})^k\bm{e}^{(0)}$. En la base de
autovectores de $\bm{H}$, la $i$-ésima componente del error evoluciona como
$(1-\eta\mu_i)^k e_i^{(0)}$. El factor de contracción global es:
\[
  \rho(\eta) = \max_i |1-\eta\mu_i|
  = \max\{|1-\eta\mu_1|,\, |1-\eta\mu_p|\}.
\]
Para convergencia se necesita $\rho(\eta) < 1$, equivalente a $0 < \eta < 2/\mu_1$.
La tasa óptima se obtiene igualando $1-\eta\mu_p = \eta\mu_1-1$, dando
$\eta^* = 2/(\mu_1+\mu_p)$, con tasa de contracción óptima:
\[
  \rho(\eta^*) = \frac{\mu_1-\mu_p}{\mu_1+\mu_p} = \frac{\kappa-1}{\kappa+1}.
\]
Por tanto $\|\bm{e}^{(k)}\|_2 \leq \rho(\eta^*)^k\|\bm{e}^{(0)}\|_2 =
\left(\frac{\kappa-1}{\kappa+1}\right)^k\|\bm{e}^{(0)}\|_2$. $\square$
\end{ejercicio}

\begin{ejercicio}
\label{ej:criticos_red}
\textbf{(Clasificación de puntos críticos en una red de una neurona).}

Sea $f(x;w,b) = \tanh(wx+b)$ con el dato $(x_1,y_1) = (1,0.5)$ y pérdida MSE.
Calcular todos los puntos críticos de $\mathcal{L}(w,b) = (\tanh(w+b)-0.5)^2$ y
clasificarlos usando el hessiano.

\medskip\noindent\textit{Solución.}
Denotemos $z = w+b$ y $t = \tanh(z)$. Entonces $\mathcal{L} = (t-0.5)^2$.

\noindent\textit{Puntos críticos.} Las derivadas parciales son:
\[
  \frac{\partial\mathcal{L}}{\partial w}
  = 2(t-0.5)\,\mathrm{sech}^2(z) \cdot x_1
  = 2(t-0.5)\,\mathrm{sech}^2(z),
\]
\[
  \frac{\partial\mathcal{L}}{\partial b}
  = 2(t-0.5)\,\mathrm{sech}^2(z).
\]
Ambas se anulan simultáneamente si $t = 0.5$ o $\mathrm{sech}^2(z) = 0$. La
segunda condición requiere $|z| \to \infty$. La primera da $z^* = \tanh^{-1}(0.5)
= \frac{1}{2}\ln 3 \approx 0.549$, y los puntos críticos forman la \emph{recta}
$w + b = z^*$ en $\mathbb{R}^2$: hay infinitos puntos críticos, todos equivalentes.

\noindent\textit{Hessiano.} En cualquier punto de la recta $w+b = z^*$, el gradiente
es $\bm{0}$ y el hessiano tiene autovalores $\{2\,\mathrm{sech}^4(z^*)\cdot 2, 0\}$:
un autovalor positivo (dirección perpendicular a la recta) y uno nulo (dirección
paralela a la recta). El índice de Morse es $0$: son mínimos locales (pero no
estrictos). La recta entera $w+b = z^*$ es un mínimo global. Esto ilustra la
simetría de una neurona: hay una manifold entera de soluciones equivalentes. $\square$
\end{ejercicio}

%───────────────────────────────────────────────────────────────────────────────
\section*{Ejercicios propuestos}
\addcontentsline{toc}{section}{Ejercicios propuestos}
%───────────────────────────────────────────────────────────────────────────────

\begin{ejercicio}
\textbf{(Gradiente y superficie de nivel).}
Sea $\mathcal{L}(\theta_1,\theta_2) = \theta_1^2 + 4\theta_2^2$.
\begin{enumerate}[label=(\alph*)]
  \item Calcule $\nabla\mathcal{L}$ y verifique la Proposición~\ref{prop:max_ascenso}
        en el punto $\bm{\theta} = (1,1)^\top$.
  \item Describa geométricamente las superficies de nivel $\mathcal{S}_c$
        para $c = 1, 4, 9$. ¿Son esferas o elipsoides?
  \item Calcule el hessiano y sus autovalores. ¿Cuál es el número de condicionamiento?
  \item Según la Proposición~\ref{prop:conv_cuad}, ¿cuántas iteraciones de descenso
        de gradiente se necesitan (aproximadamente) para reducir el error inicial
        en un factor de $10^{-3}$?
\end{enumerate}
\end{ejercicio}

\begin{ejercicio}
\textbf{(Puntos críticos de una pérdida no convexa).}
Sea $\mathcal{L}(\theta) = \theta^4 - 4\theta^2 + 1$ en $\mathbb{R}$.
\begin{enumerate}[label=(\alph*)]
  \item Calcule todos los puntos críticos y clasifíquelos (mínimos, máximos,
        inflexión) usando la segunda derivada.
  \item Trace la gráfica y verifique que la función tiene dos mínimos locales.
  \item ¿El descenso de gradiente converge al mínimo global desde cualquier
        inicialización? Identifique las cuencas de atracción de cada mínimo.
  \item Generalice: si $\mathcal{L}(\theta) = \theta^4 - 2a\theta^2 + b$ con
        $a > 0$, ¿para qué valores de $a$ tiene la función dos mínimos locales?
\end{enumerate}
\end{ejercicio}

\begin{ejercicio}
\textbf{(Condicionamiento y zigzag).}
Sea la pérdida cuadrática $\mathcal{L}(\theta_1,\theta_2) = \frac{1}{2}(\theta_1^2
+ \kappa\theta_2^2)$ con $\kappa \gg 1$.
\begin{enumerate}[label=(\alph*)]
  \item Calcule el hessiano y su número de condicionamiento.
  \item Implemente (o describa analíticamente) las primeras 10 iteraciones del
        descenso de gradiente con tasa óptima $\eta^* = 2/(\mu_1+\mu_2)$
        partiendo de $\bm{\theta}^{(0)} = (1,1)^\top$.
  \item Para $\kappa = 100$, ¿cuántas iteraciones se necesitan para que
        $\|\bm{\theta}^{(k)}\|_2 < 10^{-3}$?
  \item Compare con el número de iteraciones necesarias si se usara el hessiano
        completo (método de Newton): $\bm{\theta}^{(k+1)} = \bm{\theta}^{(k)} -
        \bm{H}^{-1}\nabla\mathcal{L}$.
\end{enumerate}
\end{ejercicio}

\begin{ejercicio}
\textbf{(No convexidad de la composición).}
Sea $f(x;\bm{\theta}) = \sigma(\theta_2\,\sigma(\theta_1 x))$ con $\sigma(z) =
\tanh(z)$ y el dato $(x_1,y_1) = (1,0.5)$.
\begin{enumerate}[label=(\alph*)]
  \item Calcule $\mathcal{L}(\theta_1,\theta_2) = (f(1;\bm{\theta})-0.5)^2$.
  \item Calcule $\nabla\mathcal{L}$ y encuentre al menos dos puntos críticos.
  \item Calcule el hessiano en cada punto crítico y determine su índice de Morse.
  \item Compare el número de puntos críticos de este modelo de dos capas con el
        de una neurona del Ejercicio Resuelto~\ref{ej:criticos_red}. ¿Cómo
        afecta la profundidad a la complejidad del paisaje?
\end{enumerate}
\end{ejercicio}

\begin{ejercicio}
\textbf{(Regularización y convexidad).}
Sea $\mathcal{L}(\bm{\theta}) = \sin(\theta_1)\cos(\theta_2) + 0.1(\theta_1^2+\theta_2^2)$.
\begin{enumerate}[label=(\alph*)]
  \item Calcule el hessiano y determine para qué valores de $(\theta_1,\theta_2)$
        el hessiano es indefinido (punto de silla).
  \item Añada regularización $\lambda\|\bm{\theta}\|_2^2$ con $\lambda > 0$.
        ¿Para qué valor mínimo de $\lambda$ la función total $\mathcal{L}(\bm{\theta})
        + \lambda\|\bm{\theta}\|_2^2$ es convexa en todo $\mathbb{R}^2$?
        (\textit{Sugerencia}: determine el mínimo autovalor de $\nabla^2\mathcal{L}$
        en todo su dominio.)
  \item Contraste con la Proposición~\ref{prop:reg_convexa}: ¿bajo qué condición
        adicional sobre $\mathcal{L}$ garantiza esa proposición la convexidad?
\end{enumerate}
\end{ejercicio}

%───────────────────────────────────────────────────────────────────────────────
\label{sec:estrella3}
\begin{estrella}{Geometría del paisaje de pérdida en la estimación inversa de reactividad}

\textbf{El sistema físico.}
La cinética puntual del reactor da la evolución de la densidad neutrónica $n(t)$
en función de la reactividad $\varrho(t)$ y las concentraciones de precursores
$C_i(t)$. El problema inverso consiste en: dado $n(t)$ medido en $T$ instantes
$\{t_k\}_{k=1}^T$, estimar el perfil de reactividad $\varrho(t)$ que lo produce.

\textbf{El paisaje de pérdida.}
Se parametriza $\varrho(t)$ con $p$ coeficientes $\bm{\theta} \in \mathbb{R}^p$
(por ejemplo, como expansión en funciones escalón o B-splines) y se define:
\[
  \mathcal{L}(\bm{\theta}) = \frac{1}{T}\sum_{k=1}^T
  \bigl(n_{\text{modelo}}(t_k;\bm{\theta}) - n_{\text{medido}}(t_k)\bigr)^2.
\]
Como $n_{\text{modelo}}(t;\bm{\theta})$ es la solución de un sistema de EDOs
no lineal (la cinética puntual), $\mathcal{L}$ es no convexa en $\bm{\theta}$.

\textbf{Geometría del paisaje.}
Los puntos críticos de $\mathcal{L}$ corresponden a perfiles de reactividad que
reproducen las mediciones de neutrones. El hessiano en un punto crítico tiene la
forma de la Proposición~\ref{prop:hess_mse}: $\bm{H} \approx \frac{2}{T}\bm{J}^\top\bm{J}$
(aproximación de Gauss-Newton), donde $\bm{J}_{kj} = \partial n(t_k)/\partial\theta_j$
es la sensibilidad de la respuesta neutrónica al $j$-ésimo parámetro de reactividad.
El condicionamiento $\kappa(\bm{J}^\top\bm{J})$ cuantifica cuánto se pueden distinguir
los distintos perfiles de reactividad a partir de las mediciones de potencia: una
dirección con $\sigma_i \approx 0$ corresponde a una perturbación de $\varrho(t)$
que produce un cambio negligible en $n(t)$ y por tanto no puede inferirse de los datos.

\textbf{Mínimos espurios.}
La no convexidad implica la posibilidad de soluciones espurias: perfiles de
reactividad $\hat{\varrho}(t)$ que reproducen $n(t)$ con pérdida baja pero son
físicamente incorrectos. En la práctica, estos mínimos locales corresponden a
oscilaciones de alta frecuencia en $\varrho(t)$ que se cancelan en la integral de la
cinética y no son detectables desde la señal de potencia. La regularización del
Capítulo~2 --- aplicada aquí como penalización $\lambda\|\bm{\theta}\|_2^2$ o
penalización de suavidad $\lambda\|\nabla\varrho\|_2^2$ --- restringe el espacio
de soluciones a perfiles suaves y elimina estos mínimos espurios.

\textbf{Conexión con el capítulo.}
Este problema ilustra los tres resultados principales del capítulo actuando
simultáneamente: (1) el hessiano en el mínimo (aproximación de Gauss-Newton)
tiene el mismo condicionamiento que la SVD de la matriz de sensibilidades $\bm{J}$;
(2) los mínimos espurios son puntos de silla en dimensiones adicionales --- se
convierten en mínimos locales solo cuando se restringe el espacio de búsqueda a
perfiles físicamente razonables; (3) la regularización convexifica la pérdida,
eliminando los mínimos espurios a costa de sesgar la solución hacia perfiles suaves.
El Capítulo~5 presentará los algoritmos iterativos para resolver este problema de
optimización; el Capítulo~7 su formulación bayesiana con incertidumbre cuantificada.
\end{estrella}

%───────────────────────────────────────────────────────────────────────────────
\section*{Ejercicios computacionales}
\addcontentsline{toc}{section}{Ejercicios computacionales}
%───────────────────────────────────────────────────────────────────────────────

\begin{ejerciciocomp}{Visualización del paisaje de pérdida en 2D}
\textbf{Objetivo:} construir y visualizar el paisaje de pérdida de un modelo
no lineal simple en $\mathbb{R}^2$ e identificar sus puntos críticos numéricamente.

\medskip\noindent\textbf{Datos:}
$N = 10$ puntos $(x_i, y_i)$ generados como $x_i = i/10$,
$y_i = \sin(2\pi x_i) + \varepsilon_i$ con $\varepsilon_i \sim \mathcal{N}(0, 0.05^2)$.

\medskip\noindent\textbf{Modelo:} familia de una capa oculta con dos neuronas y
activación $\tanh$:
\[
  f(x; \bm{\theta}) = \theta_3\,\tanh(\theta_1 x) + \theta_4\,\tanh(\theta_2 x),
\]
fijando $\theta_3 = \theta_4 = 1$ para visualizar el paisaje en el espacio
$(\theta_1, \theta_2) \in [-4, 4]^2$.

\medskip\noindent\textbf{Algoritmo:}
\begin{enumerate}
  \item Calcular $\mathcal{L}(\theta_1, \theta_2)$ en una cuadrícula de
        $100 \times 100$ puntos en $[-4,4]^2$.
  \item Graficar el paisaje como mapa de calor y curvas de nivel. Identificar
        visualmente las regiones de mínimos y puntos de silla.
  \item Calcular numéricamente $\nabla\mathcal{L}$ y $\bm{H}$ en tres puntos:
        uno en una región de mínimo aparente, uno en una cresta y uno en una
        zona plana.
  \item Para cada uno de los tres puntos, reportar los autovalores del hessiano
        y clasificar el punto (mínimo, máximo, silla, plano).
\end{enumerate}

\medskip\noindent\textbf{Verificación:}
El paisaje debe mostrar al menos dos mínimos (correspondientes a asignar
$\theta_1 \approx \pi$ y $\theta_2 \approx 3\pi$, o viceversa, por simetría de
permutación). La pérdida en el mínimo debe ser $< 0.01$. Los autovalores del
hessiano en el mínimo deben ser ambos positivos.

\medskip\noindent\textbf{Extensión:}
Repetir el análisis fijando $(\theta_1, \theta_2)$ y variando $(\theta_3, \theta_4)$.
¿El paisaje en ese subespacio es convexo? ¿Por qué?
\end{ejerciciocomp}

\begin{ejerciciocomp}{Condicionamiento del hessiano y velocidad de convergencia}
\textbf{Objetivo:} verificar experimentalmente la cota de la
Proposición~\ref{prop:conv_cuad} y observar el efecto del condicionamiento
en la trayectoria del descenso.

\medskip\noindent\textbf{Sistema:} pérdida cuadrática
$\mathcal{L}(\bm{\theta}) = \frac{1}{2}\bm{\theta}^\top\bm{H}\bm{\theta}$
con $\bm{H} = \mathrm{diag}(\mu_1, \mu_2) = \mathrm{diag}(1, 1/\kappa)$ para
$\kappa \in \{1, 10, 100, 1000\}$.

\medskip\noindent\textbf{Algoritmo:}
\begin{enumerate}
  \item Para cada $\kappa$, ejecutar descenso de gradiente con tasa óptima
        $\eta^* = 2/(\mu_1 + \mu_2)$ partiendo de $\bm{\theta}^{(0)} = (1,1)^\top$.
        Registrar $\|\bm{\theta}^{(k)}\|_2$ para $k = 0, 1, \ldots, 200$.
  \item Graficar $\log\|\bm{\theta}^{(k)}\|_2$ vs. $k$ para los cuatro valores
        de $\kappa$. Verificar que la pendiente coincide con $\log\!\left(\frac{\kappa-1}{\kappa+1}\right)$.
  \item Para $\kappa = 100$, graficar la trayectoria $(\theta_1^{(k)}, \theta_2^{(k)})$
        en el plano. Observar el zigzag.
  \item Repetir con $\eta = 2\eta^*$ (tasa demasiado grande) y registrar la
        primera iteración donde $\|\bm{\theta}^{(k)}\|_2 > \|\bm{\theta}^{(0)}\|_2$.
\end{enumerate}

\medskip\noindent\textbf{Verificación:}
Para $\kappa = 100$ y tasa óptima, el error debe caer por un factor de $10^{-3}$
en $k \approx 460$ iteraciones (calculable de la fórmula: $k \geq \frac{3\ln 10}{\ln(101/99)} \approx 460$).
Para $\kappa = 1$ (hessiano esférico), la convergencia debe ser en $k = 1$.

\medskip\noindent\textbf{Extensión:}
Implementar el método de gradiente conjugado para el mismo sistema. Comparar
el número de iteraciones necesarias para alcanzar $\|\bm{\theta}^{(k)}\|_2 < 10^{-6}$
con el descenso de gradiente. El gradiente conjugado debe converger en exactamente
$p = 2$ pasos.
\end{ejerciciocomp}

%───────────────────────────────────────────────────────────────────────────────
\begin{notasbib}
El análisis de puntos críticos de funciones diferenciables mediante el índice de
Morse fue desarrollado sistemáticamente por Marston Morse en los años 1920--1930;
su libro \textit{The Calculus of Variations in the Large} (1934) es la referencia
clásica. La teoría de Morse estudia cómo la topología del conjunto de sublevel
$\mathcal{L}_c$ cambia al variar $c$, y cómo los puntos críticos determinan esa
topología. La abundancia de puntos de silla en alta dimensión como obstáculo para
la optimización fue identificada en el contexto de redes neuronales por
Dauphin et al.~\cite{dauphin2014identifying}, quien propuso el método \textit{saddle-free
Newton} para escapar de ellos eficientemente.

La propiedad de que los mínimos locales de redes sobredeterminadas tienen valores
similares fue establecida empíricamente por Goodfellow et al.\ y analizada
teóricamente por Choromanska et al.\ (2015) usando una conexión con el modelo de
vidrios de espín de la física estadística, específicamente con el modelo de
Sherrington-Kirkpatrick. Aunque esa conexión ha sido cuestionada, la observación
empírica es robusta. La distinción entre ``muchos mínimos locales malos''
(imagen incorrecta) y ``puntos de silla prevalentes'' (imagen correcta) para
redes profundas está bien establecida~\cite{goodfellow2015qualitatively}.

El fenómeno de \textit{double descent} fue descrito formalmente por
Belkin et al.~\cite{belkin2019reconciling} y su análisis en el modelo lineal fue
completado por Hastie et al.\ (2022). La explicación en términos de regularización
implícita del SGD y sesgo hacia mínimos planos se desarrolla en el Capítulo~8.

El análisis de la tasa de convergencia del descenso de gradiente para funciones
fuertemente convexas y su dependencia del condicionamiento es estándar; la
presentación aquí sigue a Nesterov~\cite{nesterov2004introductory}. La conexión
entre el condicionamiento de la hessiana y el condicionamiento de la matriz de
diseño del Capítulo~2 está analizada en detalle en el Capítulo~8 de
Nocedal y Wright~\cite{nocedal2006numerical}, que también desarrolla los métodos
de segundo orden (Gauss-Newton, cuasi-Newton, gradiente conjugado) que el
Capítulo~5 de este libro introduce en el contexto estocástico.

La analogía entre el hessiano de la función de pérdida y la matriz de modos
normales en mecánica es explícita en la literatura de aprendizaje automático bajo
el nombre de \textit{loss surface}: el espectro del hessiano en el mínimo determina
la ``rigidez'' del modelo en el sentido de cuánto cambia la predicción ante
perturbaciones de los parámetros. Esta conexión con la mecánica cuántica de
segunda cuantización y la física estadística de sistemas desordenados es uno de
los temas activos en la teoría estadística del aprendizaje profundo.
\end{notasbib}
